\section{Direct prior art and industrial inspection systems}
\subsection{Direct academic research on coated PCBs}
The most relevant academic evidence is much narrower than the general literature on bare-PCB surface defects. The two 2021 Korean Journal of Broadcast Engineering papers by Lee et al. directly address conformal-coated PCB bubble inspection \citep{lee2021cnn,lee2021transfer}. In the earlier approach, UV images were processed to find dark-centre and bright-boundary candidates, after which CNN-based verification rejected non-bubble lookalikes. The authors' abstract reports precision 0.805, recall 0.763 and F1 0.767. As those rounded values are not algebraically consistent with the conventional harmonic-mean F1 calculated from precision and recall, the reported metric definitions and aggregation method require inspection of the complete paper before they are used as a benchmark. This discrepancy is precisely why the original evidence matters. In the later paper, VGG-style CNN transfer learning is reported to improve bubble detection to F1 0.9044. Those results cannot be directly compared without matching data, test splits and the exact positive class.

Zouhri and colleagues presented a Faster R-CNN assessment of coating bubbles at ICM 2021 and published a related 2024 journal discussion of machine-learning and edge inspection \citep{zouhri2021faster,zouhri2024ai}. Their presence further establishes the pre-existing literature on deep-learning-based coating detection; however the available bibliographic record is insufficient to evaluate the size of their dataset or reproduce their results. They remain marked for full-text follow-up.

In 2024, Choi et al. explicitly addressed UV irradiance selection for small conformal-coating bubbles, using image-level colour/brightness descriptors, feature-space distances, YOLOv8 and image tiling \citep{choi2024uv}. This is essential to the lighting and ablation design: it means our proposed illumination normalization must be positioned relative to existing UV intensity suitability approaches. Wang et al. reported a YOLOv8-Nano variant with attention and feature fusion for PCB coating bubble detection, with a 12\% relative performance improvement for smaller bubbles in the institutional abstract \citep{wang2025bubble}. That figure is \emph{source-reported and method-specific}, and cannot be treated as an absolute accuracy gain or as performance on the present boards.

More recently, Patel et al. presented a customised Faster R-CNN with ResNet-152 and enhanced proposal mechanisms for capillary and dewetting patterns on conformal-coated PCBAs \citep{patel2026coverage}. The publisher/institutional abstract reports recall of 98.53\% and 95.49\% for two defect categories, respectively. This is relevant because it moves past bubbles toward process-related coverage flaws, but recall alone says nothing about precision, false rejects, specimen independence or line deployment. A 2026 hybrid PCB inspection chapter also reports a combination of registration, component-fault learning and HSV analysis of \emph{excess conformal coating} \citep{pidurkar2026hybrid}. Its presence means the novelty of a generic classical-plus-AI pipeline must be framed very narrowly.

\subsection{Relevant adjacent work}
General PCB AOI surveys discuss sensor/illumination choices, geometric comparison, segmentation, classification and manufacturing integration \citep{ebayyeh2020aoi,zhou2023pcb,ieee2023pcb}. A 2025 PCB fault system aligns boards using RANSAC and segments defects using U-Net under several illumination conditions \citep{qub2025registration}. This is useful for designing the \emph{geometric and lighting baselines}, but its conventional PCB defects cannot validate coating-coverage decisions. Similarly, the old 1988 Hara et al. system exploited ultraviolet-excited \emph{PCB base-material fluorescence} to distinguish copper patterns \citep{hara1988fluorescent}; it is historical evidence for fluorescence-based imaging in PCB inspection, not evidence for modern coating evaluation.

Outside electronics, Han et al. used autoencoders, self-labelling, and convolutional models to identify coating boundaries on fuel-injection nozzles \citep{han2025coating}. Their fully accessible paper reports over 95\% interface localization within $\pm 6$ pixels and 7.18 images per second, along with important limitations on generalization to other sources. The concept of reducing manual annotation may transfer, but its sample geometry, surface appearance and lighting differ markedly from densely populated PCBs. Broader coating-defect surveys likewise provide method categories rather than transferable PCB inspection performance \citep{zi2026coating}.

\subsection{Commercial conformal-coating AOI systems}
\begin{table}[htbp]\centering\footnotesize
\caption{Published manufacturer descriptions. Numeric specifications are vendor claims, not independent head-to-head benchmarks.}
\begin{tabularx}{\textwidth}{p{23mm}X X}\toprule
Supplier / product & Relevant published capabilities & Engineering implication \\\midrule
Nordson FX-940UV \citep{nordsoncci} & 5 MP colour image, UV illumination, taught known-good coating/noncoating areas, optional oblique views; manufacturer describes inspection in a few seconds. & Golden reference plus rules are established practice; 5 MP can be commercially relevant for some configurations.\\
G\"opel Multi Line CCI \citep{goepelcci} & Top and side UV, CAD/coating-plan regions, claimed 9/17 $\mu$m sampling, optional multiple oblique cameras and separate thickness measurement. & Region semantics and geometric coverage plans are established; complex 3D areas demand more than one view.\\
MIRTEC GENESYS-CC \citep{mirteccci} & 50 MP camera, claimed 4 $\mu$m optics and 30 $\mu$m microbubble sensitivity; AI-based bubble assessment, additional UV and side camera geometry. & Defect size and inspection risk determine resolution; published claims need independent validation.\\
Viscom CCI \citep{viscomcci} & UV-triggered fluorescence for coating presence/evenness; distinguishes optical inspection from spot thickness measurement. & Treat fluorescence intensity as a process signal, not calibrated thickness without metrology.\\
\bottomrule\end{tabularx}
\label{tab:commercial}
\end{table}

The existence of these systems has two consequences. First, CAD-driven forbidden/required regions, UV illumination and multiple camera views cannot individually be claimed as new. Second, commercial product brochures rarely provide a publicly reproducible dataset with independent board counts and false-acceptance operating points. Academic work can add value by publishing \emph{transparent evaluation and cost-performance trade-offs} rather than simply reproducing a proprietary feature list. A credible comparison should not equate our €-level prototype estimate to a large commercial machine including conveyors, warranty, service, traceability and optional three-dimensional metrology.

\subsection{Synthesis: what has not been established by available evidence}
Direct papers demonstrate the feasibility of learning bubble and selected coverage-defect patterns, but they do not collectively establish an openly benchmarked low-cost, single-view, CPU-only system meeting a specified factory false-acceptance rate across independent coating lots. The reviewed commercial pages demonstrate use of reference plans, but do not independently establish how much they improve each specific defect class. Thus an honest research opportunity is to isolate the \emph{incremental value of reference awareness and illumination correction} under controlled tests. It remains possible that a calibrated classical thresholding algorithm will outperform a neural network for simple missing-coating or keep-out checks, while neural processing is needed only for complicated bubbles and texture defects.
